\section{Discussion and Threats to Validity}
\label{sec:discussion}

The theoretical formalization and empirical validation of the CN-OMM yield several critical implications for software architects, Site Reliability Engineers (SREs), and researchers designing cloud-native observability and AIOps platforms. This section discusses these architectural takeaways and addresses potential threats to validity.

\subsection{Architectural and FinOps Implications}
\label{subsec:implications}

\subsubsection{Why Premature AIOps Investments Fail Without Level 4 Stage-Gates}
A major practical and theoretical insight of our configurational model is explaining why enterprise investments in commercial ``AIOps'' platforms or Large Language Model (LLM) incident assistants frequently fail or stall in production deployments~\cite{amgothu_2025, soldani_2022, roy_2024_fse}. Software organizations frequently succumb to the classical \textit{stage-skipping fallacy}~\cite{humphrey_1988, lasrado_2015, vansteenbergen_2010}---attempting to leap directly from Level 2 (reactive CLI monitoring) or Level 3 (siloed tracing and unindexed logs) to Level 5 (autonomous, predictive remediation and self-healing).

Maturity model theory establishes that architectural capabilities are strictly cumulative and interdependent: higher-order optimization stages cannot function without the standardized measurement, semantic linkage, and correlation mechanisms of prerequisite levels~\cite{humphrey_1988, lasrado_2015}. In cloud-native microservices, this principle is particularly acute: as documented in recent empirical studies on cloud incident triage~\cite{wang_2023_icse, roy_2024_fse} and microservice RCA~\cite{soldani_2022}, deploying black-box machine learning models or LLMs over unaligned telemetry creates a severe ``Garbage In, Garbage Out'' failure mode characterized by hallucinations, false alarms, and upstream misattributions. Our controlled experiment with Baseline B empirically substantiates this dynamic: feeding raw inclusive trace latencies ($D(s_i)$) without exclusive self-time isolation ($T_{\text{self}}$) and W3C cross-pillar correlation ($D_2$) into an analytical engine produced $100\%$ upstream misattribution under cascading delay scenarios (Section~\ref{subsec:empirical_results}). By codifying this dependency into a non-compensatory Stage-Gate rule ($\text{Level}_{\text{eff}} = \min_k \lfloor C(D_k) \rfloor$), the CN-OMM formally prevents organizations from premature AI adoption before establishing the causal, multimodal Level 4 telemetry foundation.

\subsubsection{Elevating Probe Effect and FinOps Governance to First-Class Maturity Criteria}
Unlike vendor maturity frameworks that incentivize unbounded data ingestion~\cite{grafana_omm_2024, dynatrace_omm_2024}, Dimension $D_5$ of the CN-OMM penalizes observability architectures that violate CPU/memory stability budgets ($\Delta_{\text{CPU}} \ge 15\%$) or rely on full-text indexing of high-cardinality labels. Achieving Level 4 or Level 5 maturity explicitly requires architectural discipline---combining OpenTelemetry collector batching, tail-based anomaly sampling, and index-free object storage (such as Grafana Loki and Tempo)---so that diagnostic entropy reduction is achieved sustainably.

\subsection{Threats to Validity}
\label{subsec:threats}

Following Wohlin et al.~\cite{wohlin_2012}, we analyze four categories of validity threats:
\begin{itemize}[leftmargin=*, nosep]
    \item \textbf{Construct Validity:} To avoid subjective maturity self-reporting (criticized by Proen\c{c}a~\cite{proenca_2016_rcis} and Wendler~\cite{wendler_2012}), all 15 sub-components of the CN-OMM are operationalized through declarative Quantitative Capability Indicators ($\text{QCI}_{k.j}$) and objective system metrics ($MTTD, MTTDiag, Acc, \Delta_{\text{CPU}}$).
    \item \textbf{Internal Validity:} In the Type-3 chaos evaluation, background noise and warm-up transients were mitigated by enforcing a 60-second steady-state stabilization baseline before each fault injection, flushing telemetry buffers between runs, and pairing all $N=30$ trials across identical workload seeds and fault intensities.
    \item \textbf{External Validity:} While the CNCF OpenTelemetry Astronomy Shop is a widely recognized, realistic polyglot benchmark (14 microservices across 7 languages), ultra-large-scale clusters comprising thousands of microservices or serverless/FaaS topologies may exhibit additional control-plane bottlenecks. Nevertheless, the mathematical formulation of $\mathcal{O}(t)$, $T_{\text{self}}(s_i)$, and the 15 QCIs are agnostic to cluster size.
    \item \textbf{Conclusion Validity:} We avoided parametric assumptions on skewed latency distributions by applying non-parametric Wilcoxon signed-rank tests, Vargha-Delaney $\hat{A}_{12}$ effect sizes, and exact binomial/Fisher/McNemar tests with explicit 95\% confidence intervals.
\end{itemize}
